% ===== §3 The Modes Redrawn =====
\section{The Modes Redrawn: A Morphology of Happiness from the Standpoint of the Generative Relation}
\label{sec:modes}

\noindent
The modes sketched, pre-theoretically, in the delineation of the phenomenon are here \emph{redrawn}---reorganised in the language grounded by \xref{sec:generativity} (the generative relation, self-continuation, the modes of lack and fullness, the two ways of living the gap). The delineation was a sketch; this chapter is the redrawing after theoretical grounding---the same phenomenon passing a second time. In keeping with \xref{sub:gen-selflimit}, the chapter does \emph{not} argue that happiness ``must'' take some form; it describes the modes in which one and the same fullness emerges under varying conditions, and the list is \emph{open}. The morphology is descriptive and open, not an essential determination. Unlike the delineation, this chapter may use the full vocabulary, since \xref{sec:generativity} has grounded it.

\subsection{From sketch to redrawing}
\label{sub:modes-sketch}

The sketch is recovered (uneventful, full, shared, faintly moving, without want; with waiting as a near relative). The gain of the redrawing: the sketch could say only ``it is still and full''; after \xref{sec:generativity} one can say \emph{why}---still, because generation does not settle itself (no event, no remainder); full, because the gap, lived in the mode of fullness, drives a generation ``between''; faintly moving, because that continual generation is itself the movement (the ``slight liveliness'' the sketch felt is now nameable as the persistence of generation).

\subsection{Mode one---abiding: a being-with already in continuance}
\label{sub:modes-abiding}

The pole of ``doing nothing, watching the rain together'': the relation has fallen into a self-sustaining continuance, with no need to be oriented toward anything not yet arrived. Redrawn: the fullness of a generative relation in its self-continuation, the gap operating in the \emph{mode of fullness}---the gap is still there (hence generation, hence the faint movement), but the subject is not driven by lack (hence no want, no event, the silence). In the dynamical illustration (\xref{sec:dynamics}) this mode answers to a stable co-rotating trajectory---\emph{but only as a side-image, never as a definition}; no formal object is unfolded here, only the interface is marked.

\subsection{Mode two---waiting: a proceeding toward a common rhythm}
\label{sub:modes-waiting}

The pole of ``slowing for one another, watching the evening together'': the fullness is \emph{under way}, oriented toward what is not yet arrived (the evening descending, the other catching up).

Redrawn (load-bearing; with \xref{sub:gen-twomodes} and \xref{sec:coexistence}):
\begin{itemize}
  \item Waiting is \emph{not stillness, not lack}. The genuine (sweet) waiting is full of movement oriented toward the other---a continual adjustment of oneself, a voluntary restraint of one's own pace. Beneath the apparent stillness is an actively sustained orientation.
  \item Waiting therefore confirms, rather than refutes, ``fullness comes from continual generation'': it is a \emph{proceeding phase} of the mode of fullness---the open space the gap holds is lived as a ``moving toward together,'' not as a ``deficit to be filled.''
  \item The line against ``forced empty waiting'' (sketched in the delineation, now given its frame): sweet waiting is the gap lived in the mode of fullness (voluntary, oriented, shared); fretful empty waiting is the gap lived in the mode of lack (compelled, out of control, oriented to a deficit). \emph{One and the same ``waiting,'' two modes}---one further layer of the distinction of \xref{sub:gen-twomodes}, of Paper~XVIII's pathological/healthy coupling, of the two neural systems of \xref{sub:neuro-waiting}.
\end{itemize}
In the dynamical illustration, this is more like an approach trajectory, a negotiation not yet phase-locked, than a closed orbit---and it is stated plainly: to force it into any single form (such as a limit cycle) would be to repeat the error of ``swallowing an ill-fitting phenomenon with a cherished form.'' Waiting is the living proof that happiness is exhausted by no single form.

\subsection{Mode three---\ldots (left open): the chapter's reticence}
\label{sub:modes-open}

There are other modes, emerging historically and materially, which the paper does not foreclose and cannot enumerate. \emph{This ellipsis is itself the chapter's load-bearing element}: it positively enacts the ``leaving open of the essence of happiness'' of \xref{sub:gen-selflimit}---the openness of the morphology is not an oversight but a structural honouring of ``no finite form is entitled to exhaust happiness.'' One thread may be laid here as a footnote without elevating it into another pinned-down cell: a mode of ``never repeating yet always bounded,'' the happiness of a creative life---each day different yet always within a certain form of life. It is left as a thread for some future paper, not developed; the ``strange attractor'' figure that might tempt one belongs in this footnote, not in the body.

\subsection{The methodological discipline of the morphology}
\label{sub:modes-discipline}

The three modes (abiding / waiting / \ldots) are not the exhaustive cells of a classificatory system but several identifications within an open family. What they share is not a fixed structure but the same \emph{generation---operating in the mode of fullness of the gap, rhythmed by self-continuation, occurring ``between.''} And this ``sharing'' is itself a saying in motion, open to revision under later conditions (per the voice-discipline of \xref{sub:gen-selflimit}). The transition: the modes redrawn, the chapters that follow turn to \emph{illustration}---the dynamical (\xref{sec:dynamics}), the neural (\xref{sec:neuro})---casting analysable, visible side-light on these modes. Illustration is subordinate to the philosophical morphology established here; it does not define it in reverse.
